% !TEX root = ../main.tex
% =====================================================================
%  §5  Mapping The UV Theory To The IR Theory   （原 PDF pp. 37--57）
%  公式编号 (5.1)--(5.92)，逐条保留原文编号。
% =====================================================================

\section{从紫外理论到红外理论的映射}

\begin{keytakeaways}
\item \textbf{物理图像：} 拓扑不变关联函数的 UV 路径积分只依赖于度规的 Q-精确形变，因此可把度规整体放大（$g_{\mu\nu}\to t g_{\mu\nu}$，$t\to\infty$）后在低能有效理论（LEET）中计算；但紧流形上必须对所有量子真空求和，因为隧穿会把它们混合起来。
\item \textbf{数学工具：} 库仑分支的商结构 $(\mathfrak{t}\otimes\CC-\Delta)/W$、塞伯格--威滕椭圆曲线族与其全纯微分 $\lambda_{SW}$、周期与预势、模形式（$\Gamma_{0}(4)$）。低能数据由“全纯地变化的一族阿贝尔簇 $A\to B^{*}$ 加一个中心荷函数”给出。
\item \textbf{核心结论：} 在一般四维流形 $X$ 上
$Z^{\xi}_{IR}(p,\Sigma)=Z_{\text{Coulomb}}+Z_{\text{Higgs}}$：前者是 $u$ 平面上的积分（带引力耦合 $E(u)^{\chi}P(u)^{\sigma}$ 与通量扇区求和相位），后者局部化为塞伯格--威滕方程的解模空间上的积分。UV 与 IR 的算子映射并非环同态，其偏差正是接触项 $\Sigma^{2}T(u)$。
\end{keytakeaways}

\subsection{研究真空结构的动机}

\begin{keytakeaways}
\item \textbf{物理图像：} 拓扑不变关联函数的 UV 路径积分在度规整体放大 $g\to tg$（$t\to\infty$）下不变（这是 $Q$-精确形变），而 $t\to\infty$ 等价于把能量标度到零，故必可在 LEET 中计算；紧流形上则必须对所有量子真空求和（体积有限时隧穿不被压低）。
\item \textbf{数学工具：} LEET 是靶空间为量子真空模空间的非线性 $\sigma$ 模型，一般不可重整化；但拓扑场论只需在真空附近做到单圈阶，故仍可计算。
\item \textbf{核心结论：} 还需 UV$\to$IR 的 $Q$-闭算子对应 (5.3)(5.4)。$Q$-闭算子构成（分次）交换环，而 RG 映射\emph{不是}环同态——(5.1) 中的 $+\cdots$ 即来源于此，它最终由接触项给出。
\end{keytakeaways}

对于拓扑不变的关联函数，配分函数 $Z^{\xi}_{DW}(p,\Sigma)$——它由 UV 路径积分定义——应当可以用低能有效作用量来计算：
\begin{equation}
Z^{\xi}_{DW}(p,\Sigma):=\left\langle \ee^{p\mathcal{O}(\wp)+\mathcal{O}(\Sigma)}\right\rangle_{UV}
=Z^{\xi}_{IR}(p,\Sigma):=\left\langle \ee^{p\mathcal{O}_{IR}(\wp)+\mathcal{O}_{IR}(\Sigma)+\cdots}\right\rangle_{IR}
\tag{5.1}
\end{equation}
理由在于我们可以把度规整体放大：将
\begin{equation}
g_{\mu\nu}\to t\,g_{\mu\nu}
\tag{5.2}
\end{equation}
替换，并取极限 $t\to+\infty$。

一方面，改变 $t$ 在路径积分中是一个 $Q$-精确的改变：它不能改变积分值。另一方面，从物理观点看，我们是在把长度拉伸到无穷，相应地也就把能量标度到零。也就是说，我们研究的是真空之上无穷小尺度的动力学。因此，配分函数必定可以在低能有效理论中计算。（这正是 LEET 的定义！）

我们的目标是尽可能把 (5.1) 写得显式。我们将看到，它是克龙海默--姆罗夫卡结构定理 \cite{ref21} 与威滕猜想 \cite{ref55} 两者的推广。

\medskip
\noindent\textbf{评注：}
\begin{enumerate}[label=\arabic*., leftmargin=2.1em]
\item 我们将发现，$\RR^{3}\times\RR$ 上的 $\mathcal{N}=2$ 理论没有唯一的量子真空，而是有一个量子真空的模空间。在紧流形 $X$ 上计算该理论的路径积分时，必须对所有真空积分。原因是量子涨落在紧流形上导致这些真空之间的隧穿。隧穿振幅随 $X$ 的体积指数压低，并在 $X$ 非紧时趋于零；在此情形下，我们把无穷远处取定的某个真空作为边界条件的一部分。\footnote{〔原稿边注〕数学家觉得这一点极难把握。此处应作进一步解释。}
\item LEET 涉及一个非线性 $\sigma$ 模型，其靶空间就是方才提到的量子真空模空间。因此在 LEET 中对场作积分，就包含了对这些量子真空的积分。
\item 一般而言 LEET 不可重整化，描述 $\mathcal{N}=2$ 场论的 LEET 也不例外。因此关于这种路径积分的意义还有许多话要说。所幸在拓扑场论中，只需在某个真空附近把微扰论做到至多单圈阶，因此这些不可重整化场论的路径积分实际上是可算的。
\item 我们还需要 UV 与 IR 理论之间 $Q$-闭算子的精确对应：
\begin{equation}
\mathcal{O}(\wp)\to\mathcal{O}_{IR}(\wp)
\tag{5.3}
\end{equation}
\begin{equation}
\mathcal{O}(\Sigma)\to\mathcal{O}_{IR}(\Sigma)
\tag{5.4}
\end{equation}
粗略地说，这是算子在 RG 流下到 IR 的映射。尽管 $Q$-闭算子构成一个（分次）交换环，RG 映射却\emph{不是}环同态。这正是 (5.1) 右端指数中 $+\cdots$ 的来源，下面会把它补全。此事与所谓“接触项”的课题有关。
\end{enumerate}

\subsection{经典真空与自发对称破缺}

\begin{keytakeaways}
\item \textbf{物理图像：} 经典真空是零能态，由 $[\phi,\phi^{*}]=0$、$\mu(M)=0$、$\phi\cdot q=\phi^{*}\cdot\tilde q^{*}=0$ 刻画；解集 $\mathcal{V}$ 是 $G$-空间，真空模空间为 $\mathcal{V}/G$。
\item \textbf{数学工具：} 半单正则的 $\phi$ 给出库仑分支 $(\mathfrak{t}\otimes\CC-\Delta)/W$；$\phi=0$ 给出经典希格斯分支 $W/\!/\!/\!/G$（超凯勒商）；其余为混合分支，故 $\mathcal{M}^{\text{Classical}}$ 是分层空间。
\item \textbf{核心结论：} $SU(2)$ 纯规范理论中 $u^{\text{classical}}=\Tr\phi^{2}=2a^{2}$ 参数化 $\CC-\{0\}$：$u=0$ 处 $W^{\pm}$ 变轻、稳定化子跳为 $SU(2)$。经典真空的退化正是下一小节要论证“不被量子修正解除”的对象。
\end{keytakeaways}

现在研究 Minkowski 度规下 $\mathbb{M}^{1,3}$ 上的理论。“真空”是 Hilbert 空间中能量最低的态。在相对论性理论中，它们将具有相对论不变性，从而给出我们在 Wick 转动后的欧氏理论中做微扰论的展开点。

具有李群 $G$ 与 HM 表示 $W=R\oplus\bar R$ 的一般经典拉氏量，其 Hamilton 量是非负项之和。理论的经典真空全都具有零能量。把各类导数项置零会使规范场为零，并使标量场成为与位置无关的常数。我们把常数值记作 $\phi,q,\tilde q$。它们不可与同名的场混淆——后者是 $\mathbb{M}^{1,3}$ 上的函数。（这些常数值应被理解为“经典真空期望值”。）

令势能项为零，得到\footnote{〔原注〕我们把质量参数取为零。}
\begin{equation}
[\phi,\phi^{*}]=0
\tag{5.5}
\end{equation}
\begin{equation}
\mu(M)=0
\tag{5.6}
\end{equation}
\begin{equation}
\phi\cdot q=0\qquad \phi^{*}\cdot\tilde q^{*}=0
\tag{5.7}
\end{equation}
满足 (5.5)--(5.7) 的配对 $(\phi,M)$ 的解集 $\mathcal{V}$ 是一个 $G$-空间，而经典真空的模空间按定义为
\begin{equation}
\mathcal{M}^{\text{Classical}}:=\mathcal{V}/G
\tag{5.8}
\end{equation}

让我们略微描述一下空间 $\mathcal{M}^{\text{Classical}}$。方程 (5.5) 说明 $\phi$ 是半单的。因此它可以共轭到极大环面 $\mathfrak{t}\otimes\CC$ 中。一般而言，它的 vev 会自发破坏规范对称性：
\begin{equation}
G\to\Stab(\phi)
\tag{5.9}
\end{equation}
当 $\phi$ 既半单又正则时，稳定化子恰为嘉当环面的正规化子，未破缺的规范群恰为极大环面 $T$，其李代数为 $\mathfrak{t}$。此时方程 (5.6) 与 (5.7) 给出 $M=0$。于是 $\mathcal{M}^{\text{Classical}}$ 的一个子集可以等同于
\begin{equation}
\mathcal{M}^{\text{classical}}_{\text{Coulomb}}=(\mathfrak{t}\otimes\CC-\Delta)/W
\tag{5.10}
\end{equation}
其中 $W$ 是外尔群，$\Delta$ 是半单但非正则元素的子集。这称为真空的\emph{库仑分支}。称之为“库仑分支”，是因为未破缺的阿贝尔规范对称性意味着 LEET 中存在无质量的阿贝尔规范场，它们会对理论中带电的粒子施加电磁力；特别地，两个静止的带电粒子将感受到长程库仑力。

% [Tutor Note: (5.10) 的商结构值得细看。(i) 分子 t⊗C−Δ 是“半单且正则”的位形集合：正则性要求中心化子恰为 Cartan 子代数 t（而非更大），这是把 φ 共轭到 t⊗C 后“未破缺群 = T”的唯一条件；Δ 恰是 φ 的稳定化子严格大于 T（即某些根 α 满足 ⟨α,φ⟩=0，使 W± 无质量）的轨迹，也就是 Weyl 腔室的墙壁。因此 (5.10) 真正是“正则元素的 Weyl 商”。(ii) 但作者把 Δ 仅定义为“半单但非正则元素”，这与后文 §5.7.1 所用的判别式 Δ 同号，需注意二者是不同对象：此处的 Δ 是复超平面并集 ∪_α{⟨α,φ⟩=0}，而 (5.62)(5.63) 的 Δ 是 SW 曲线的判别式。读者应按上下文区分。(iii) 商去 Weyl 群后得到的正是 Weyl 腔室的内部；在 SU(2) 情形它就是一个复平面去掉原点（(5.16)），而“正则”这一限制在 u 平面上的像正是 u≠0。故 (5.10) 的奇点集与“规范群增强”的点不是一回事：Δ 的像在经典理论里就是量子修正可能产生新分支的地方。]

当 $\phi\in\Delta$ 时，真空方程仍有解，但部分非阿贝尔规范对称性得以恢复，同时 $M$ 也可以取非零值。其极端情形是点 $\phi=0$：若无 HM，则规范对称性完全恢复；若有 HM，则存在一个真空分支，它恰是 $W$ 的超凯勒商。这称为\emph{经典希格斯分支}：
\begin{equation}
\mathcal{M}^{\text{classical}}_{\text{Higgs}}=W/\!/\!/\!/G
\tag{5.11}
\end{equation}
称之为希格斯分支，是因为若 $M$ 处于一般位置，则 (5.7) 的解迫使 $\phi=0$，从而完全破坏连续规范对称性。于是规范对称性被希格斯化。

一般而言 $\mathcal{M}^{\text{Classical}}$ 是一个复杂的分层空间
\begin{equation}
\mathcal{M}^{\text{Classical}}=\coprod_{\alpha}\mathcal{M}^{\text{Classical}}_{\alpha}
=\mathcal{M}^{\text{Classical}}_{\text{Coulomb}}\amalg\cdots\amalg\mathcal{M}^{\text{Classical}}_{\text{Higgs}}
\tag{5.12}
\end{equation}
各分量的闭包以复杂的方式相交。其余分量满足 $\phi\in\Delta-\{0\}$，故半单但非正则，且 $M$ 也非零。这些称为“混合分支”，文献中对它们的考察并不比上述分支少多少。

\noindent\textbf{例：} 我们的主要例子是无物质超多重态的 $G=SU(2)$。经典理论中真空能量为 $V=\Tr([\phi^{*},\phi])^{2}$。它被正规矩阵极小化，因此可以通过规范变换把 $\phi$ 化为
\begin{equation}
\phi=
\begin{pmatrix}
a & 0\\
0 & -a
\end{pmatrix}
\tag{5.13}
\end{equation}
的形式，且在外尔变换 $a\to-a$ 下唯一。经典真空由规范不变参数
\begin{equation}
u^{\text{classical}}=\Tr\phi^{2}=2a^{2}
\tag{5.14}
\end{equation}
参数化。$u^{\text{classical}}$ 可以取复平面中的任意值；反之，$u$ 的一个取值唯一确定 $\RR^{4}$ 上理论的一个经典真空（原文此处把这个参数误写作 $v$）。在 $u^{\text{classical}}$ 平面上每个 $a\neq0$ 的点，规范群自发破缺：
\begin{equation}
SU(2)\to U(1)
\tag{5.15}
\end{equation}
而 $W^{\pm}$ 玻色子（即与对角生成元之外的生成元相对应的规范场）的质量正比于 $|a|$。经典地，在 $u^{\text{classical}}=0$ 处这些规范玻色子变为无质量，稳定化子跳变为 $\Stab(\phi)=SU(2)$，完整的 $SU(2)$ 对称性恢复。因此对 $SU(2)$：
\begin{equation}
\mathcal{M}^{\text{classical}}_{\text{Coulomb}}=\CC-\{0\}
\tag{5.16}
\end{equation}

\begin{figure}[H]
  \centering
  \figplaceholder{〔原文此处于第 40 页留有图位，PDF 中未插入实际图形〕}
  \caption{经典库仑分支。对无超多重态的 $G=SU(2)$，经典库仑分支为 $\CC-\{0\}$，即图中所示的复 $u^{\text{classical}}$ 平面。在 $u^{\text{classical}}=0$ 处非阿贝尔的 $SU(2)$ 规范对称性未破缺。}
  \label{fig:uclassical}
\end{figure}
% [原 PDF 第 40 页]

\subsection{量子真空与 LEET}

\begin{keytakeaways}
\item \textbf{物理图像：} 弱耦合下经典真空的退化不被解除，但真空附近的低能涨落描述会被量子效应改变；强耦合还可产生新的纯量子真空分支。
\item \textbf{论证链：} 绕经典真空做微扰论 $\to$ 非微扰效应由绕瞬子解展开定义 $\to$ 该量子理论有 LEET $\to$ $\mathcal{N}=2$ 超荷是守恒流积分，无超对称反常，故 RG 流不破缺超对称 $\to$ LEET 由 $\mathcal{N}=2$ 超对称作用量描述。
\item \textbf{核心结论：} 正因为 $\mathcal{N}=2$ 作用量不允许 $\mathcal{N}=2\to\mathcal{N}=0$ 的自发破缺（这与 $\mathcal{N}=1$ 中可诱导超势/FI 参数形成对照），原真空必为精确量子真空；经典真空退化因此不被解除，只是其邻域的低能描述会被强量子效应重构。
\end{keytakeaways}

要点如下：

\smallskip
\noindent\emph{在弱耦合的量子理论中，经典真空的退化不会被解除；但量子效应会改变对该真空附近低能涨落的描述。此外，由于强量子效应，可能出现纯量子力学的真空分支。}
\smallskip

“经典真空退化不被解除”的论证如下。选定一个经典真空，并以在该真空附近的展开来定义微扰论意义下的理论。非微扰效应则通过在瞬子解附近定义微扰论来定义。这个量子理论有一个 LEET 描述。由于 $\mathcal{N}=2$ 超荷是守恒流的时间积分，只要没有超对称反常，超对称就不会被到 IR 的 RG 流破坏；而完全没有证据表明存在超对称反常。因此我们假定量子理论中确实没有超对称反常，于是 LEET 将由一个 $\mathcal{N}=2$ 超对称有效作用量描述。

然而：$\mathcal{N}=2$ 超对称作用量不允许自发超对称破缺 $\mathcal{N}=2\to\mathcal{N}=0$，所以原来的真空必定已经是一个精确的量子真空。

（这里 $\mathcal{N}=2$ 超对称是关键的。在 $\mathcal{N}=1$ 超对称作用量中，可以诱导出超势和（或）FI 参数，二者中任何一个都能破坏超对称。）

如何描述与某个经典真空态相联系的理论的 LEET？只要没有新的量子态随着真空参数的变化而变为无质量，我们就可以用在经典真空附近展开所得的低能自由度来描述 LEET。

为实施这一点，先为简单起见假定我们处理的是紧单群 $G$ 的单个 VM。于是在理论的通常量子化中，必须引入一个标度 $\Lambda$，在该标度处耦合变为强耦合。引入裸耦合常数\footnote{〔原稿边注〕要赋予这一切意义，还需给出 YM 项的正规化方式。}
\begin{equation}
\tau_{0}=\frac{4\pi\ii}{g_{0}^{2}}+\frac{\theta_{0}}{2\pi}
\tag{5.17}
\end{equation}
以及微扰 Feynman 图上的截断 $\Lambda_{0}$，并通过固定如下组合
\begin{equation}
\Lambda^{2h^{\vee}}:=\Lambda_{0}^{2h^{\vee}}\ee^{2\pi\ii\tau_{0}}
\tag{5.18}
\end{equation}
来定义以 $\Lambda$ 为参数的量子理论，同时令 $\Lambda_{0}\to\infty$ 且 $g_{0}\to0$。\footnote{〔原注〕有物质时此式须作修改：把 $-2h^{\vee}$ 换成上文标有 [待补] 处定义的 $\beta$。}

% [Tutor Note: (5.18) 是本节唯一的“重整化群定义”。它并不定义 $\Lambda$ 本身，而是把组合 $\Lambda^{2h^{\vee}}=\Lambda_{0}^{2h^{\vee}}e^{2\pi i\tau_{0}}$ 在 $\Lambda_{0}\to\infty$、$g_{0}\to0$ 的极限下固定住；因此 $\Lambda$ 是由匹配得出的动力学生成标度，只有其 $2h^{\vee}$ 次幂才是 RG 不变的（单圈意义下）。指数恰为 $2h^{\vee}$，是因为一圈 $\beta$ 函数的系数就是 (4.76) 中的 $-2h^{\vee}+C_{2}(R)$（纯 SYM 时即 $-2h^{\vee}$）：$\tau$ 的虚部按 $1/g^{2}$ 跑动，故 $\Lambda^{2h^{\vee}}$ 不跑。脚注 19 提醒有物质时须把 $-2h^{\vee}$ 换成该 $\beta$；这也解释了为何瞬子修正 (5.32) 以 $\Lambda^{2kh^{\vee}}$（而非 $\Lambda^{2k}$）为展开参数。]

现在理论中有了标度 $\Lambda$；在自然度规下，当 $u$ 远离非正则轨迹时，理论实际上是弱耦合的。对 $SU(2)$ 而言判据是 $|u|\gg|\Lambda^{2}|$，此时可以可靠地计算经典库仑分支上某个真空附近涨落的低能谱。

把这一推理应用于我们的主要例子（无 HM 的 $G=SU(2)$），可以说至少当 $|u|\gg|\Lambda^{2}|$ 时，量子真空 $|\Omega(u)\rangle$ 可由
\begin{equation}
2u:=\langle\Omega(u)|\mathcal{O}|\Omega(u)\rangle
=\langle\Omega(u)|\frac{\Tr_{\mathbf{2}}\phi^{2}}{8\pi^{2}}|\Omega(u)\rangle
\tag{5.19}
\end{equation}
参数化——至少在我们确知量子理论中没有额外无质量态时如此。等式左端的因子 $2$ 是为了便于与数学结果作比较。

在其 1994 年春的突破性论文中，塞伯格与威滕提出——依据诸如半经典单值性、以及用超势 $W=m\Tr\phi^{2}$ 扰动所得的 $\mathcal{N}=1$ SYM 的质量间隙与禁闭等物理理由——在量子理论中库仑分支的真空并不是 $\CC-\{0\}$，而是整个 $u$ 平面。也就是说，对 $u$ 平面上\emph{每}一点，规范群都自发破缺为 $SU(2)\to U(1)$。然而在 $u=\pm\Lambda^{2}$ 处确有特殊现象发生：在这些点上，某些与“BPS 粒子”相联系的量子态变为无质量，并出现新的“量子希格斯分支”。对 $\RR^{4}$ 上的真空而言，这只是 $W=\CC\oplus\CC$ 的 HK 商，是一个单点。在一般四维流形 $X$ 上，我们将得到与塞伯格--威滕方程的解相联系的真空。\footnote{〔原注〕在一般四维流形上我们滥用了“真空”一词。不过在 $X=\Sigma_{3}\times\RR$ 上确实有真空的概念：它们应与 $\Sigma_{3}$ 上塞伯格--威滕方程解的模空间相关。在量子理论中我们应考虑该空间上的波函数，所得的 Hilbert 空间 $\mathcal{H}(\Sigma_{3})$ 应与某一类塞伯格--威滕--Floer 同调相关。}

于是 $SU(2)$ 理论的真空可以图示如下：\footnote{〔原稿边注〕距第 5 节开始已过 11 分钟。}

\begin{figure}[H]
  \centering
  \figplaceholder{〔原文此处于第 41 页留有图位，PDF 中未插入实际图形〕}
  \caption{复 $u$ 平面上的两个奇点。新无质量 HM 出现的两个特殊点位于 $u=\pm\Lambda^{2}$。在经典极限 $\Lambda^{2}\to0$ 下这两点并合，从而恢复到真空的经典图像。量子理论中存在重要的极限次序问题：在 $u$ 平面上任何 $u\neq\pm\Lambda^{2}$ 的点，LEET 由 $U(1)$ 的 $\mathcal{N}=2$ 阿贝尔规范理论描述；但在一个小的非零能标上，$u=\pm\Lambda^{2}$ 的邻域只有用与一个电荷为 $1$ 的超多重态耦合的 $U(1)$ $\mathcal{N}=2$ 阿贝尔规范理论才能准确描述。}
  \label{fig:uplane}
\end{figure}
% [原 PDF 第 41 页]

\subsubsection{描述 $\RR^{4}$ 上库仑分支涨落的一般 $\mathcal{N}=2$ 超对称 LEET}

只含阿贝尔矢量多重态、规范群为阿贝尔群 $T$（李代数 $\mathfrak{t}$）的一般 $\mathcal{N}=2$ LEET，由一个非线性 $\sigma$ 模型描述，其靶空间为
\begin{equation}
\mathcal{B}=(\mathfrak{t}\otimes\CC)/W
\tag{5.20}
\end{equation}
该描述将在某些余维一子簇的补集 $\mathcal{B}^{*}$ 上成立；越过这些子簇，谱中就会出现阿贝尔矢量多重态所描述之外的新无质量态。

$\mathcal{N}=2$ 超对称的一条重要定理是：映入 $\mathcal{B}^{*}$ 的场的 LEET 一般作用量，由 $\mathcal{B}^{*}$ 上一族阿贝尔簇的数据
\begin{equation}
\pi:\mathcal{A}\to\mathcal{B}^{*}
\tag{5.21}
\end{equation}
连同全纯变化的 $\mathcal{N}=2$ \emph{中心荷函数} $Z$ 所决定。$Z$ 应视为全空间上的全纯函数：
\begin{equation}
Z:\Gamma\to\CC
\tag{5.22}
\end{equation}
它在纤维上线性。这里 $\Gamma$ 是一个格的局部系统
\begin{equation}
\Gamma\to\mathcal{B}^{*}
\tag{5.23}
\end{equation}
其纤维为 $\Gamma_{u}:=H_{1}(\mathcal{A}_{u};\ZZ)$。这些格带有反对称的整值“Dirac--Schwinger--Zwanziger”配对。该局部系统通常具有非平凡的单值性。

在任一点 $u\in\mathcal{B}^{*}$ 的某个无单值性邻域 $\mathcal{U}$ 中，我们可以取 $\Gamma$ 的纤维的一个极大拉格朗日分解：
\begin{equation}
\Gamma_{u}=\Gamma_{u}^{\text{electric}}\oplus\Gamma_{u}^{\text{magnetic}}
\qquad u\in\mathcal{U}
\tag{5.24}
\end{equation}
这样的选择称为一个\emph{对偶标架的选取}。

式 (5.24) 中给 $\Gamma_{u}$ 的子格标上“电”与“磁”，是因为格 $\Gamma_{u}$ 的物理意义正是未破缺阿贝尔规范对称性下的电磁荷格；那个整值反对称配对则解释为电磁荷的 Dirac 配对。这可由下文所述拉氏量连同阿贝尔 $S$-对偶得到。

与任一给定的对偶标架相联系，有一组局部坐标
\begin{equation}
a^{I}=Z(\alpha^{I})
\tag{5.25}
\end{equation}
其中 $I$ 遍历 $\Gamma_{u}$ 的秩。此外，$\mathcal{N}=2$ 超对称还给出一个局部定义的全纯函数 $\mathcal{F}$，称为\emph{预势}，使得
\begin{equation}
\frac{\partial\mathcal{F}}{\partial a^{I}}=Z(\beta_{I})
\tag{5.26}
\end{equation}
这些量通常就定义为 $a_{D,I}$。最后，阿贝尔簇 $\mathcal{A}_{u}$ 相对于该对偶标架的周期矩阵为
\begin{equation}
\tau_{IJ}:=\frac{\partial a_{D,I}}{\partial a^{J}}
=\frac{\partial^{2}\mathcal{F}}{\partial a^{I}\partial a^{J}}
\tag{5.27}
\end{equation}

% [Tutor Note: (5.27) 把三个表面上不同的对象等同起来：周期矩阵 $\tau_{IJ}$（阿贝尔簇 $\mathcal{A}_{u}$ 相对所选对偶标架的周期）、磁坐标对电坐标的导数 $\partial a_{D,I}/\partial a^{J}$、以及预势的二阶偏导。其全部内容在于 $\mathcal{F}$ 的\emph{全纯性}：因此 $\tau_{IJ}$ 对 $a$ 全纯，而不只是实解析。这一点是 LEET 的强约束——$\tau$ 不能任意给定，必须是某个全纯预势的 Hessian。相应地 LEET 标量场的 Kähler 势可写成 $K=\mathrm{Im}\big(\bar a^{I}\partial\mathcal{F}/\partial a^{I}\big)$，而规范动能正比于 $\mathrm{Im}\,\tau_{IJ}$，故物理要求 $\mathrm{Im}\,\tau>0$（正定性）。还要注意指标的协变位置：$a^{I}$ 为上标（电），$a_{D,I}$ 为下标（磁），$\tau_{IJ}$ 对称全纯；后文 §5.4 的 $a_{D}=\partial\mathcal{F}/\partial a$ 就是秩一情形。]

现在，用这些数据，LEET 的作用量为
\begin{equation}
S_{IR}=\int \ii\left(\tau_{IJ}F^{-,I}F^{-,J}+\bar\tau_{IJ}F^{+,I}F^{+,J}\right)
+\mathrm{Im}\,\tau_{IJ}\,\dd a^{I}\ast \dd\bar a^{J}+\cdots
\tag{5.28}
\end{equation}
其中 $+\cdots$ 中的各项都涉及费米子以及由 $\tau_{IJ}$ 及其导数定义的非线性相互作用。

对紧单群 $G$ 的矢量多重态理论，在库仑分支上某个弱耦合真空附近定义其量子理论，一般结果是：存在一族特殊的对偶标架，其特殊坐标为 $a\in\mathfrak{t}\otimes\CC$，而“预势”为
\begin{equation}
\mathcal{F}=\mathcal{F}^{\text{classical}}+\mathcal{F}^{1\text{-loop}}+\mathcal{F}^{\text{instanton}}
\tag{5.29}
\end{equation}
\begin{equation}
\mathcal{F}^{\text{classical}}+\mathcal{F}^{1\text{-loop}}
=\frac{\ii}{4\pi}\sum_{\alpha\in\Delta^{+}}\langle\alpha,a\rangle^{2}
\log\frac{\langle\alpha,a\rangle^{2}}{2\Lambda^{2}}
\tag{5.30}
\end{equation}
其中 $\Delta^{+}$ 是 $\mathfrak{g}$ 的一组正根。为了把上式分离为经典与单圈贡献，我们利用式 (5.18)，得到
\begin{equation}
\mathcal{F}^{\text{classical}}=\text{const.}\,\tau_{0}(a,a)
\tag{5.31}
\end{equation}
非微扰修正具有如下形式\footnote{〔原稿边注〕把常数算出来。}
\begin{equation}
\mathcal{F}^{\text{instanton}}
=\sum_{k=1}^{\infty}\mathcal{F}_{k}\big(\langle\alpha,a\rangle^{-1}\big)\Lambda^{2kh^{\vee}}
\tag{5.32}
\end{equation}
其中 $\mathcal{F}_{k}$ 是 $|\Delta^{+}|$ 个变量的 $2kh^{\vee}-2$ 次齐次多项式。

\medskip
\noindent\textbf{评注：}
\begin{enumerate}[label=\arabic*., leftmargin=2.1em]
\item 利用 \cite{ref39,ref40} 中发展的技巧，涅克拉索夫及其合作者已从第一性原理出发，对所有经典群、以及带有各种半单群与物质的推广，算出了这些齐次多项式 $\mathcal{F}_{k}$ \cite{ref43,ref44,ref45}。
\item 应当强调，当 (5.28) 被用作某个 UV 理论的 LEET 时，它只是导数展开——即按 $E/\Lambda$（$E$ 为典型能标）展开——中的领头项。然而对于本文所讨论的拓扑问题，只有导数展开的领头项是重要的：更高阶项全都是 $Q$-精确的，在任何情形下都不会影响低能极限。
\end{enumerate}

\subsection{塞伯格与威滕关于 $G=SU(2)$ 的 LEET}

\begin{keytakeaways}
\item \textbf{物理图像：} 在 $u\neq\pm\Lambda^{2}$ 处 LEET 只含一个 $U(1)$ 矢量多重态 $(a,A,D,\eta,\chi,\psi)$，即从 $\RR^{4}$ 到量子库仑分支的 $\sigma$ 模型；$\mathcal{N}=2$ 只要求 $\tau(a)$ 全纯，给出 $\tau(a)$ 就等于指定了低能理论。
\item \textbf{数学工具：} 拓扑扭曲后的拉氏量 (5.35) 把所有 $\tau$ 依赖统一写成 $\tau(a)$ 及其对 $\bar a$ 的导数；绕大 $u$ 的单值性由威滕效应给出 (5.34) 的矩阵。
\item \textbf{核心结论：} 确定 LEET 归结为两步：算出 $\tau(a)$，并说明 $\tau$ 与 $u$ 的关系；这正是下一小节用塞伯格--威滕椭圆曲线族所要解决的问题。
\end{keytakeaways}

回到塞伯格与威滕最初分析的无 HM 的 $G=SU(2)$ 情形，LEET 的描述如下：

在 $u=\pm\Lambda^{2}$ 之外的任一真空附近展开时，LEET 中的场包含单个 $U(1)$ 矢量多重态：
\begin{equation}
a,\,A,\,D,\,\eta,\,\chi,\,\psi
\tag{5.33}
\end{equation}
取值于李代数 $\mathfrak{u}(1)$。因此 $a$ 是 $\RR^{4}$ 上的复值标量场。概念上它应被理解为定义了一个非线性 $\sigma$ 模型，即映射 $\RR^{4}\to\mathcal{M}^{\text{Quantum}}_{\text{Coulomb}}$。$A$ 是阿贝尔规范场，其余类推。

\medskip
\noindent\textbf{评注：}
\begin{enumerate}[label=\arabic*., leftmargin=2.1em]
\item 物理学家总是用同一个 $a$ 同时描述我们作展开所围绕的量子 VEV，以及描述该 VEV 长波涨落的量子场。
\item 如前所强调，为了写出 LEET 必须选取一个对偶标架。其他对偶标架中的描述，可通过在 $U(1)$ 规范理论上作电磁对偶得到——这在路径积分中是可以显式实施的 \cite{ref46,ref56}。尽管 $u$ 很大处的局部系统有单值性，它可由威滕效应很好地理解。因此从物理可知，某个生成元的作用为
\begin{equation}
\begin{pmatrix}
a\\
a_{d}
\end{pmatrix}
\to
\begin{pmatrix}
-1 & 0\\
4 & -1
\end{pmatrix}
\begin{pmatrix}
a\\
a_{d}
\end{pmatrix}
\tag{5.34}
\end{equation}
特别地，$a$ 在差一个符号的意义下是良定义的。（事实上，由上面给出的预势表达式可以立即导出这一点。）\footnote{〔原稿边注〕核对因子 2。}
\end{enumerate}

拉氏量为
\begin{align}
\mathcal{L}_{IR,vm}\sim{}&
\ii\left(\bar\tau F^{+}F^{+}+\tau F^{-}F^{-}\right)\notag\\
&+\mathrm{Im}\,\tau\,\dd a\ast \dd\bar a+\mathrm{Im}\,\tau\,D\ast D\notag\\
&+\tau\psi\ast \dd\eta+\bar\tau\eta\,\dd\ast\psi+\tau\psi\,\dd\chi-\bar\tau\chi\,(\dd\psi)
\tag{5.35}\\
&+\ii\frac{\dd\tau}{\dd\bar a}\eta\chi\left(D+F^{+}\right)+\cdots\notag
\end{align}

% [Tutor Note: (5.35) 的符号与因子约定是本节最容易出错的地方，作者自己在页边写下“Check factor of 2”。逐项看：第一行取 $\bar\tau F^{+}F^{+}+\tau F^{-}F^{-}$，即自对偶部分配 $\bar\tau$、反自对偶部分配 $\tau$；这与 (5.28) 的一般形式 $\tau_{IJ}F^{-,I}F^{-,J}+\bar\tau_{IJ}F^{+,I}F^{+,J}$ 一致（只是把 $\tau$ 写在 $F^{-}$ 之前）。若按常见文献写成 $\ii(\tau F^{+}F^{+}+\bar\tau F^{-}F^{-})$，则 $\tau$ 与 $\bar\tau$ 的角色整体互换，$\mathrm{Im}\,\tau$ 的符号约定随之改变。第二处要点是末项 $\ii(d\tau/d\bar a)\eta\chi(D+F^{+})$：$a$ 全纯、$\bar\tau$ 反全纯，故对 $\tau$ 求导只能对 $\bar a$；写 $\dd\tau/\dd\bar a$ 而非 $\dd\tau/\dd a$ 是扭曲形式中的关键。第三，常数因子：$a$ 的归一化、(5.34) 中的矩阵元 $4$、以及 (5.19) 中定义的 $2u:=\langle\mathcal{O}\rangle$，是与 (5.38)/(5.41) 的模函数展开逐项匹配后才确定的；作者标注的 factor of 2 提醒读者不要把它与 $e^{2\pi i\tau}$ 与 $e^{\pi i\tau}$ 的约定混为一谈。]

\medskip
\noindent\textbf{评注：}
\begin{enumerate}[label=\arabic*., leftmargin=2.1em]
\item 这里给出的是我们所需要的那种拓扑扭曲形式；$+\cdots$ 中包含其他复杂的相互作用项，我们并不需要（但在非单连通流形上它们会变得相关）。
\item 就 $\mathcal{N}=2$ 超对称的约束而言，$\tau(a)$ 可以是 $a$ 的任意全纯函数。给出 $\tau(a)$ 就等于指定了拉氏量。因此，为确定低能理论，我们需要：
\begin{enumerate}[label=\alph*.), leftmargin=1.8em]
\item 计算 $\tau(a)$；
\item 说明 $\tau(a)$ 如何与 $u$ 相联系。
\end{enumerate}
\end{enumerate}

塞伯格与威滕对这一问题的解答如下：（我们并不试图论证这一解答，只是以适合本文目的的方式简洁陈述结果。）

\begin{enumerate}[label=\arabic*., leftmargin=2.1em]
\item 考虑如下椭圆曲线族：\footnote{〔原注〕就计算而言，这实际上并不是该曲线族最方便的表示；见下文 §5.4.1。}
\begin{equation}
E_{u}:\qquad y^{2}=(x-u)(x-\Lambda^{2})(x+\Lambda^{2})
\tag{5.36}
\end{equation}
注意判别式轨迹显然位于 $u=\pm\Lambda^{2}$。
\item 给这些曲线配上亚纯微分
\begin{equation}
\lambda_{SW}:=\frac{\dd x}{y}(x-u)
\tag{5.37}
\end{equation}
于是中心荷函数由 $\lambda_{SW}$ 的周期给出。注意该微分在 $x=\infty$ 处有极点；全纯微分是 $\omega=\dd x/y$。
\item 族 (5.36) 所对应的局部系统确实具有由 (5.34) 中生成元的转置所给出的单值性。我们选取 $H_{1}(E_{u})$ 的一组拉格朗日同调基 $A,B$，使得 $A$ 在该生成元的两倍作用下不变，并\emph{定义}
\begin{equation}
a=\oint_{A}\lambda_{SW}.
\tag{5.38}
\end{equation}
\item 于是 $\tau$ 就是 $E_{u}$ 相对于该同调基的周期——这给出函数 $\tau(a)$，并且我们（至少隐含地）说明了 $a$ 的（vev）如何与真空 $u$ 相联系，其关系由 (5.38) 给出。注意，正如 LEET 讨论中的惯例，我们对 $\RR^{4}$ 上的场（例如 $a(x^{\mu})$）与其在 $\RR^{4}$ 上的真空期望值使用同一个记号 $a$。
\item 当然，我们选取了一组同调基，但有效理论并不依赖于这一选取，因为拉格朗日同调基的改变对应于阿贝尔理论上的一个电磁对偶变换。在 $u$ 平面的某些区域存在优先的选取：在 $u\to\infty$ 的区域，我们应选取同调基使得
\begin{equation}
\mathrm{Im}\,\tau(u)\to\infty
\tag{5.39}
\end{equation}
因为 $\tau=\frac{4\pi\ii}{e^{2}}+\frac{\theta}{2\pi}$。
\end{enumerate}

% [Tutor Note: (5.38) 是把几何与电荷格接起来的关键约定。$\lambda_{SW}$ 是第二类亚纯微分（在 $x=\infty$ 有极点），其周期 $(a,a_{D})=(\oint_{A}\lambda_{SW},\oint_{B}\lambda_{SW})$ 给出中心荷 $Z=\gamma_{e}\,a+\gamma_{m}\,a_{D}$；因此 $A$-周期对应电荷、$B$-周期对应磁荷。由 $\lambda_{SW}$ 的极点留数可知 $\lambda_{SW}$ 的周期在模变换下有平移，这正是 (5.34) 中矩阵下三角形状的来源。注意 (5.37) 中的因子 $(x-u)$ 正是使 $\lambda_{SW}$ 在 $u\to\infty$ 时具有正确经典极限 $\lambda_{SW}\sim a$ 的选取；若省去它，$\omega=\dd x/y$ 只给出 $u$ 的解析函数而非 $a$。至于“$A$ 在生成元的两倍作用下不变”，这一条件把同调基唯一确定到 $\Gamma(2)$；而 §5.4.1 改用 $\Gamma_{0}(4)$ 的框架后，$A$-周期与 $B$-周期的定义在模曲线上的表述会更自然（(5.41)(5.42)）。]

这一关于真空结构的解答导致一个非常引人注目的现象：局部系统 $H_{1}(E_{u};\ZZ)$ 在判别式轨迹 $u=\pm\Lambda^{2}$ 周围具有单值性，而纤维化 $E\to\CC$ 正是在那里变得奇异。在这些点上，$\lambda_{SW}$ 的某个周期趋于零。这将与某些非常重要的物理相联系。\footnote{〔原稿边注〕应当把局部系统描述得更清楚！选定割线、基，并给出绕 3 条显式回路的矩阵。}

\begin{figure}[H]
  \centering
  \figplaceholder{〔原文此处于第 45 页留有图位，PDF 中未插入实际图形〕}
  \caption{复 $u$ 平面上的环面丛在两点处退化。}
  \label{fig:torifibration}
\end{figure}
% [原 PDF 第 45 页]

\subsubsection{另一种描述}

塞伯格--威滕曲线族可以有多种呈现方式。上面给出的是 $\Gamma(2)$ 的模区域上的一族曲线。用定义在 $\Gamma^{0}(4)$ 的模区域上的曲线，有等价的描述：
\begin{equation}
y^{2}=x^{2}(x-u)+\frac{\Lambda^{4}}{4}x
\tag{5.40}
\end{equation}
这略好一些，因为它更容易推广到包含 HM 表示的曲线。判别式轨迹仍在 $u=\pm\Lambda^{2}$。

利用相对于无穷远处优先基的 $\tau$，$u$ 与 $\tau$ 的关系可以用模函数非常漂亮地表达出来：
\begin{equation}
u=\frac{1}{2}\frac{\vartheta_{2}^{4}+\vartheta_{3}^{4}}{(\vartheta_{2}\vartheta_{3})^{2}}
=\frac{1}{8q^{1/4}}\left(1+20q^{1/2}-62q+216q^{3/2}+\cdots\right)
\tag{5.41}
\end{equation}
对偶标架并未被完全固定，因为 $B$-圈的选择并不唯一。单值性给出 $\tau\to\tau+4$，同时保持 $u$ 不变，这是理所当然的。

类似地，$a(u)$ 与 $a_{d}(u)$ 也可以用模函数非常显式地表达。例如
\begin{equation}
a(u)=\frac{1}{6}\left(\frac{2E_{2}+\vartheta_{2}^{4}+\vartheta_{3}^{4}}{\vartheta_{2}\vartheta_{3}}\right)
=\frac{1}{4}q^{-1/8}\left(1+O(q^{1/2})\right)
\tag{5.42}
\end{equation}
$u$ 平面可以等同于 $\Gamma^{0}(4)$ 的基本域，其中 $\tau=0$ 与 $\tau=2$ 两个尖点分别对应于磁单极点与双荷子点。

\begin{figure}[H]
  \centering
  \figplaceholder{〔原文此处于第 46 页留有图位，PDF 中未插入实际图形。原文图内标注：$\Gamma^{0}(4)$ 的基本域；半经典区域：$\mathcal{F}\cup T\cdot\mathcal{F}\cup T^{2}\cdot\mathcal{F}\cup T^{3}\cdot\mathcal{F}$；磁单极区域：$S\cdot\mathcal{F}$；双荷子区域：$T^{2}S\cdot\mathcal{F}$〕}
  \caption{$\Gamma^{0}(4)$ 的基本域及其分解：半经典区域为 $\mathcal{F}\cup T\cdot\mathcal{F}\cup T^{2}\cdot\mathcal{F}\cup T^{3}\cdot\mathcal{F}$，磁单极区域为 $S\cdot\mathcal{F}$，双荷子区域为 $T^{2}S\cdot\mathcal{F}$。}
  \label{fig:gammazero4}
\end{figure}
% [原 PDF 第 46 页]

% [Tutor Note: (5.41)(5.42) 的 $q$-展开可用来校核归一化约定：由 (5.41) 得 $u\sim q^{-1/4}/8$，由 (5.42) 得 $a\sim q^{-1/8}/4$，于是 $a^{2}\sim q^{-1/4}/16=u/2$，即 $u=2a^{2}$，与 (5.14) 的经典关系一致——这正是 (5.19) 中把 $2u$ 而非 $u$ 定义为 $\langle\mathcal{O}\rangle$ 的配套约定。另需注意符号：原文用上指标记法 $\Gamma^{0}(4)$（而非 $\Gamma_{0}(4)$），它与 $\Gamma(2)$ 通过 $G(\tau):=F(\tau/2)$ 相联系；正文中“$\Gamma_{0}(4)$”与图题中的 $\Gamma^{0}(4)$ 是同一对象的不同排版，翻译时统一取原文图像的上指标形式。此外 $E_{2}$ 不是模形式而是准模形式，其模变换会产生常数项——这正是后文引入 $\hat E_{2}$ (5.92) 的原因。]

一般而言，若 $F(\tau)$ 对 $\Gamma(2)$ 是模的，则 $G(\tau):=F(\tau/2)$ 对 $\Gamma^{0}(4)$ 是模的。因此这两族椭圆曲线是同源的。

\medskip
\noindent\textbf{评注：}
\begin{enumerate}[label=\arabic*., leftmargin=2.1em]
\item 除纯 $SU(2)$ 规范理论这一情形外，我们并无把 $u$ 平面等同于某条模曲线的便利。因此最好设法使用那些能推广到任意库仑分支的论证。不过我们仍会利用把 $u$ 平面等同于 $\Gamma^{0}(4)$ 模区域所带来的某些简化。
\item 事实上，基于 class $S$ 理论的现代观点给出了写出塞伯格--威滕曲线与微分的相当不同（且通常更优越）的方式。但这里我们采用较早的表述，因为 $u$ 平面上的积分正是在那一语言中完成的。
\end{enumerate}

\subsubsection{一个重要的对称性}

对紧单群 $G$ 的单个 VM 的 $\mathcal{N}=2$ 理论，量子理论有一个重要的整体对称性。它是 $U(1)_{R}$ 对称群的有限子群。一般而言，瞬子效应把该 $U(1)_{R}$ 明显破缺为有限群 $\ZZ_{2h^{\vee}}$。随后，库仑分支上的真空期望值又把
\begin{equation}
\ZZ_{2h^{\vee}}\to\ZZ_{2}
\tag{5.43}
\end{equation}
进一步破缺。$\ZZ_{2}$ 不被破缺，因为它等价于空间中的 $2\pi$ 旋转。

用 $\phi$ 来说，我们可以取 $\ZZ_{2h^{\vee}}$ 的一个生成元使其作用为
\begin{equation}
\phi\to\omega\phi
\tag{5.44}
\end{equation}
其中 $\omega=e^{2\pi\ii/2h^{\vee}}$。由此可知库仑分支上的规范不变参数满足
\begin{equation}
u_{k}\to\omega^{k}u_{k}
\tag{5.45}
\end{equation}
这里 $u_{k}\sim\Tr\phi^{k}$ 是 0-可观测（量）。

特别地，对 $G=SU(2)$ 我们有 $u\to-u$ 的对称性。它把 $u=+\Lambda^{2}$ 处的奇点映到 $u=-\Lambda^{2}$ 处的奇点。这非常有用，可使我们在计算 $Z^{\xi}_{DW}$ 的希格斯分支贡献时工作量减半。

\subsubsection{一般的塞伯格--威滕曲线}

带“夸克”物质超多重态（表示取 $R=\mathbf{2}^{\oplus N_{f}}$，$0\le N_{f}\le4$，用式 (4.65) 的记号）的 $SU(2)$ 理论的 LEET，已由塞伯格与威滕 \cite{ref47,ref46} 解出。他们的讨论已被许多其他物理学家推广到大量理论；有用的综述见 \cite{ref51}。

正如刚才提到的，对大量理论，塞伯格--威滕曲线已由涅克拉索夫等人用瞬子微积分导出。还存在一大类理论——“class $S$ 理论”——它们与希钦系统密切相关；对这些理论，SW 曲线就是希钦底空间上的谱曲线，可以直接写出。

然而应当强调，对于任意的 $\mathcal{N}=2$ 场论，如何写出 SW 曲线族与 SW 微分至今仍然未知。

\subsection{BPS 态}

\begin{keytakeaways}
\item \textbf{物理图像：} LEET (5.35) 在 $u=\pm\Lambda^{2}$ 处变为奇异：按 $\Gamma^{0}(4)$ 的基本域看，两个尖点 $\tau=0,2$ 处 $\mathrm{Im}\,\tau\to0$，是强耦合极限，作用量中许多项发散。这表明 (5.35) 在 $u=\pm\Lambda^{2}$ 的小邻域内是错的：该 LEET 是积分掉有质量模式得到的，若有粒子在 $u_{*}$ 处变轻，就必须把它们“积分进来”。
\item \textbf{数学工具：} 固定无穷远处真空 $u$ 的 Hilbert 空间按电磁荷格分次 $\mathcal{H}_{u}=\oplus_{\gamma\in\Gamma_{u}}\mathcal{H}_{u,\gamma}$；中心荷 $Z_{\gamma}(u)$ 在 $\mathcal{N}=2$ 超庞加莱代数中按 c-数作用，Bogomolnyi 界给出 $H\ge|Z_{\gamma}(u)|$，BPS 态按定义饱和该界；精确质量公式为 $M=\text{const.}\,|\oint_{\gamma}\lambda_{SW}|$。
\item \textbf{核心结论：} $A$-圈周期给电荷、$B$-圈周期给磁荷；$u=\Lambda^{2}$ 处 $B$-周期 $a_{D}=\oint_{B}\lambda_{SW}$ 消失，对应的正是一个无质量的磁荷 $1$ 磁单极子，$u=-\Lambda^{2}$ 处则是 $B+A$ 圈对应的双荷子。它们必须以超多重态身份重新进入 LEET，从而迫使对偶标架变换 (5.53)/(5.56)。
\end{keytakeaways}

LEET (5.35) 在 $u=\pm\Lambda^{2}$ 处变为奇异。若用 $\Gamma^{0}(4)$ 的基本域来理解 $u$ 平面，这一点尤为明显：在两个尖点 $\tau=0,2$ 处有 $\mathrm{Im}\,\tau\to0$，即强耦合极限，作用量中许多项将发散。这是一个信号，表明以 (5.35) 为作用量的 LEET 在 $u=\pm\Lambda^{2}$ 的小邻域内是\emph{错的}。

LEET 失效的原因在于它是由积分掉有质量模式得到的。相当一般地，假定理论（作为真空参数 $u$ 的函数）的谱使得某些粒子在 $u$ 平面上某点 $u_{*}$ 变为无质量。那么为了得到非奇异的 LEET，对应于这些粒子的场必须被“积分进来”。也就是说，需要找到一个新的 LEET，它既包含描述 $u$ 平面上涨落的 $U(1)$ VM，又包含新的轻场，使得一旦把这些轻场“积分掉”就回到原来的 LEET（在低于这些轻场质量的能标上）。所谓“积分掉”，就是对那些轻场作路径积分，但不对 $U(1)$ VM 的场作，并保留低能展开的领头项。当人们试图积分掉无质量粒子时，Feynman 图会出现奇点，于是所谓的 LEET 中就会出现奇点。

由此可以断定，在 $u=\pm\Lambda^{2}$ 处，理论谱中必定存在某种无质量粒子，它不同于描述库仑真空附近低能涨落的无质量 $\mathcal{N}=2$ Maxwell 场多重态。

注意这里的极限次序问题：对任一固定的 $u\neq\pm\Lambda^{2}$，只要“典型的”场位形不延伸到 $u=\pm\Lambda^{2}$，理论就由 (5.35) 很好地描述。然而，若要得到一个对取值于某个包含（比如）$u=\Lambda^{2}$ 的开邻域 $\mathcal{U}_{\Lambda^{2}}$ 的场位形都成立的描述，就必须把在 $u=\pm\Lambda^{2}$ 处变轻的场“积分进来”。

一般而言，即便在相对好处理的 $d=4$ $\mathcal{N}=2$ 理论中，求 Hamilton 量作为 $u$ 的函数的谱也是不可能的。所幸理论的 Hilbert 空间中存在一个子空间，其 Hamilton 量的谱是精确已知的。这就是所谓的 \emph{BPS 谱}。简言之，BPS 谱可以如下定义：$d=4$ $\mathcal{N}=2$ QFT 在 $\RR^{3}$ 上、在 $|\vec x|\to\infty$ 处真空为 $u$ 的 Hilbert 空间记作 $\mathcal{H}_{u}$，它按未破缺阿贝尔规范对称性下的电磁荷格 $\Gamma_{u}$ 分次：
\begin{equation}
\mathcal{H}_{u}=\oplus_{\gamma\in\Gamma_{u}}\mathcal{H}_{u,\gamma}
\tag{5.46}
\end{equation}
在超选择扇区 $\mathcal{H}_{u,\gamma}$ 中，$\mathcal{N}=2$ 超庞加莱代数里的 $\mathcal{N}=2$ 中心荷\emph{算子} $Z$ 按 c-数 $Z_{\gamma}(u)$ 作用。此外，通过考虑超荷的适当线性组合的 Hermite 平方并利用超对称代数，可以证明限制到 $\mathcal{H}_{\gamma,u}$ 上的 Hamilton 量有下界 $|Z_{\gamma}(u)|$。这个界称为“Bogomolnyi 界”。于是 BPS 谱按定义为电荷 $\gamma$ 的 BPS 态对 $\gamma$ 的求和，其中 BPS 态由
\begin{equation}
\mathcal{H}^{BPS}_{u,\gamma}:=\left\{\Psi\in\mathcal{H}_{u,\gamma}\,\big|\,H\Psi=|Z_{\gamma}(u)|\Psi\right\}.
\tag{5.47}
\end{equation}
定义。

要找出 BPS 谱，必须解决两个问题：首先必须给出 $\mathcal{N}=2$ 中心荷 $Z_{\gamma}(u)$ 的显式公式；其次必须真正确定 Hilbert 空间中哪些态饱和 Bogomolnyi 界。第一步由塞伯格--威滕曲线给出，因为 $Z_{\gamma}(u)$ 是由某个亚纯微分形式（即塞伯格--威滕微分）沿曲线各圈的周期给出的。

对纯 $SU(2)$ 理论，精确的 BPS 谱已在塞伯格--威滕的原始论文中提出。\footnote{〔原注〕如今可以先用半经典谱求解、再用孔采维奇--索伊贝尔曼穿墙公式来证明它。}既然知道了 $\lambda$ 的周期，我们就有了精确的质量公式：
\begin{equation}
M=\text{const.}\,\Big|\oint_{\gamma}\lambda_{SW}\Big|
\tag{5.48}
\end{equation}
它给出电磁荷为 $\gamma$ 的 BPS 粒子的质量。

现在回忆，电磁荷格与纤维 $\Gamma_{u}$ 等同：电磁荷等同于同调类 $\gamma\in H_{1}(E_{u};\ZZ)$。特别地，$A$-圈对应电荷，对偶的 $B$-圈对应磁荷。

% [Tutor Note: $A$/$B$ 周期与电荷格的对应，以及“周期为零 $\Rightarrow$ 粒子无质量”的逻辑值得单独点明。由 (5.48)，一个 BPS 态的质量正比于其电磁荷 $\gamma\in H_{1}(E_{u};\ZZ)$ 上的周期；故若某条回路的周期在 $u_{*}$ 处趋于零，该荷扇区的 BPS 态质量便趋于零。$A$-圈对应电荷、$B$-圈对应磁荷，而 $B$-周期 $a_{D}=\oint_{B}\lambda_{SW}$ 在 $u=\Lambda^{2}$ 处消失，于是磁荷为 $1$ 的态在那里变轻——这正就是塞伯格--威滕的“无质量磁单极子”。这是绕开“磁单极子作为经典孤子解”的另一条路：弱耦合区（$|u|\gg|\Lambda^{2}|$）中磁单极子确实可由单极子模空间的集体坐标量子化得到（数学上是单极子模空间上 Dirac 型算子的 $L^{2}$ 核），而周期论证自动给出它在强耦合端的命运。$B+A$ 圈在 $u=-\Lambda^{2}$ 处消失，给出双荷子。另需注意：本页两处“$|u|\ll|\Lambda^{2}|$”按上下文（“弱耦合区”“我们在该处使用的光子场 $A$”）疑为 $|u|\gg|\Lambda^{2}|$ 之误，译文照原文保留。]

在弱耦合区 $|u|\ll|\Lambda^{2}|$，可以通过对磁单极子模空间作集体坐标量子化来证明磁荷 BPS 粒子的存在。这是一个有趣的故事，在数学上涉及单极子模空间上 Dirac 型算子的 $L^{2}$ 核的研究。\footnote{〔原注〕这是一个历史悠久、贡献者众多的长故事；近期的讨论见 \cite{ref41} 及其所引文献。}特别地，磁荷为 $1$ 的磁单极子带有一个作为 $B$-圈的荷 $\gamma$，并且这个 $B$-圈具有如下性质：周期
\begin{equation}
a_{d}:=\oint_{B}\lambda_{SW}
\tag{5.49}
\end{equation}
在 $u=\Lambda^{2}$ 处为零。这就是塞伯格与威滕著名的“无质量磁单极子”。

在单极子超多重态正在变轻的邻域 $\mathcal{U}_{\Lambda^{2}}$ 中，它应作为超多重态场被纳入 LEET。由于该粒子带磁荷，相应的场 $M$ 通过
\begin{equation}
\int\left|(\dd+A_{d})M\right|^{2}+\cdots
\tag{5.50}
\end{equation}
与 $U(1)$ 光子耦合，其中 $A_{d}$ 是磁对偶的 $U(1)$ 光子。它与我们在 $|u|\ll|\Lambda^{2}|$ 处所用的光子场 $A$ 由
\begin{equation}
F(A)=\ast F(A_{d})
\tag{5.51}
\end{equation}
相联系。因此，为了写出局域的 LEET，我们必须作适当的对偶标架变换。$A_{d}$ 的超对称伙伴是周期
\begin{equation}
a_{d}=\oint_{B}\lambda_{SW}
\tag{5.52}
\end{equation}
电荷格与磁荷格互换，于是
\begin{equation}
\tau_{d}=-1/\tau
\tag{5.53}
\end{equation}
（这里用 $\Gamma^{0}(4)$ 的基本域来描述非常方便。）

现在回忆理论具有 $u\to-u$ 的对称性。必定存在一个类似的粒子在 $u=-\Lambda^{2}$ 处变为无质量，它是与圈 $B+A$ 相对应的“双荷子”。一般而言，“双荷子”是同时带电、磁荷的粒子。半经典分析表明，除磁单极子外还存在一个磁荷为 $1$、电荷为 $1$ 的双荷子，它对应于圈 $B+A$，并在 $u=-\Lambda^{2}$ 处变为无质量。\footnote{〔原注〕更确切地说，由于无穷远处电荷格的单值性，磁荷为一、电荷为偶数的双荷子可以在 $u=\Lambda^{2}$ 处变为无质量，而磁荷为一、电荷为奇数的双荷子可以在 $u=-\Lambda^{2}$ 处变为无质量。}

\medskip
\noindent\textbf{评注：}
\begin{enumerate}[label=\arabic*., leftmargin=2.1em]
\item 在一般的 $d=4$、$\mathcal{N}=2$ 场论中，我们其实期望在 $\mathcal{N}=2$ 场论里能精确求解\emph{两}件事。除 LEET 之外，我们还期望能求解“BPS 谱”。这些是（低能阿贝尔规范理论）某个固定荷扇区中最轻的稳定粒子。SW 同样求出了纯 $SU(2)$ 理论的 BPS 谱；此后在理解许多其他理论的该谱方面有很大进展，但一般解仍未获得。这个迷人的课题本身就需要一门课程……
\end{enumerate}

\subsection{$u=\pm\Lambda^{2}$ 附近的低能理论}

\begin{keytakeaways}
\item \textbf{物理图像：} (5.35) 在 $u=\pm\Lambda^{2}$ 处失效，但在小邻域 $\mathcal{U}_{\pm\Lambda^{2}}$ 内修补方式明确：把在那些点变轻的 BPS 粒子（以超多重态表示）对应的场“积分进来”。邻域的大小取决于描述 LEET 所用的能量截断。
\item \textbf{数学工具：} 由于 $A$ 与 $A_{d}$ 只有非局域的关联，要写局域拉氏量就必须整块地改用对偶矢量多重态 $(a_{d},A_{d},\dots)$；对偶标架变换在 $u=+\Lambda^{2}$ 处是 $\tau_{d}=-1/\tau$ (5.53)，在 $u=-\Lambda^{2}$ 处是 $\tau_{d}=-1/(\tau+2)$ (5.56)。
\item \textbf{核心结论：} 这样得到的 $S_{IR,\mathcal{U}_{\Lambda^{2}}}$ 在拓扑扭曲后取上同调场论的标准形式，从而局域化到塞伯格--威滕方程 (5.57)；这正是 $Z_{\text{Higgs}}$ 的来源，而 $u\to-u$ 对称性使 $-\Lambda^{2}$ 处的计算可由 $+\Lambda^{2}$ 处的结果得到。
\end{keytakeaways}

总结上一小节：尽管描述 (5.35) 在 $u=\pm\Lambda^{2}$ 处失效，但如何在这些点附近的小邻域 $\mathcal{U}_{\pm\Lambda^{2}}$ 内修正低能描述是清楚的。（邻域的大小取决于我们用来描述 LEET 的能量截断。）我们需要把在 $u=\pm\Lambda^{2}$ 处变为无质量的粒子所对应的场“积分进来”。这些粒子最终发现是处于超多重态表示中的 BPS 粒子。

于是，举例来说，在 $u=\Lambda^{2}$ 处变轻的粒子是一个超多重态，它在弱耦合时是磁荷为 $1$、电荷为 $0$ 的磁单极子。描述该粒子的场，在 $u=\Lambda^{2}$ 附近的小邻域内就只是一个超多重态，它在 $U(1)$ VM $(a,\eta,\chi,\psi,A)$ 的电磁对偶下具有电荷 $+1$。

于是 $u=+\Lambda^{2}$ 邻域内的 LEET 由一对场多重态描述：
\begin{enumerate}[label=\arabic*., leftmargin=2.1em]
\item 一个 $U(1)$ VM，其标量场为 $a_{d}$、矢量场为 $A_{d}$，即 $A$ 的电磁对偶；
\item 一个 HM，由对偶电磁规范群 $U(1)_{D}$ 的电荷 $+1$ 表示定义。
\end{enumerate}

看出为什么需要使用对偶规范场的一种方式是：单极子标量场 $M$ 在大 $|u|$ 处的对偶标架中带磁荷 $+1$，因此当把它们作为场纳入拉氏量时，其动能项必须取
\begin{equation}
|DM|^{2}=|\dd M+A_{d}M|^{2}
\tag{5.54}
\end{equation}
的形式，其中 $A_{d}$ 是电磁对偶场。由于 $A$ 与 $A_{d}$ 是非局域相关的，而我们希望使用局域的拉氏量，所以在写出矢量多重态对作用量的贡献时，应当使用对偶的矢量多重态场 $(a_{d},A_{d},\dots)$。\footnote{〔原稿边注〕下面这一段前面已经讨论过，是否删去？}

这一 LEET 的作用量记作 $S_{IR,\mathcal{U}_{\Lambda^{2}}}$，以强调它应被用于 $u=\Lambda^{2}$ 附近的小区域。为得到它，取对偶标架由 (5.53) 定义的作用量 (5.35)，加上标准 HM 拉氏量，并作拓扑扭曲；细节可参见 \cite{ref24}。于是该 LEET 的作用量具有如下形式：
\begin{equation}
S_{IR,\mathcal{U}_{\Lambda^{2}}}
=S_{IR,VM}(a_{d},A_{d},\eta_{d},\chi_{d},\psi_{d})
+S_{HM}(M=q\oplus\tilde q^{*})
\tag{5.55}
\end{equation}
在拓扑扭曲之后——即耦合到背景 $SU(2)_{R}$ 对称性联络并令 $\omega^{+}=\omega_{R}$——上述作用量取上同调场论的标准形式，局域化到 SW 方程上。

在 $u=-\Lambda^{2}$ 处有非常类似的故事。这里圈 $B+A$ 的周期为零。相关的模变换为
\begin{equation}
\tau_{d}=-\frac{1}{\tau}+2
\tag{5.56}
\end{equation}
而使理论无质量的是磁荷为 $1$、电荷为 $1$ 的双荷子（在 $W^{\pm}$ 玻色子带电荷 $\pm2$ 的单位制中）。\footnote{〔原稿边注〕必须说明如何平凡化局部系统，这才能讲得通！！}所得理论的作用量记作 $S_{IR,\mathcal{U}_{-\Lambda^{2}}}$。实践中，一旦把 $u=+\Lambda^{2}$ 处的计算做完，便可利用 §5.4.2 的有用对称性得到答案。\footnote{〔原稿边注〕距 §5.3.1 开始已过 33 分钟。}

\subsection{把扭曲的 LEET 放到一般四维流形 $X$ 上}

\begin{keytakeaways}
\item \textbf{物理图像：} $X$ 的体积有限，真空之间的隧穿振幅不被压低，因此必须对所有真空求和。这一求和的一部分，就是计算描述 $u$ 平面上涨落的 LEET 的配分函数；但在 $u=\pm\Lambda^{2}$ 的小邻域内必须改用 (5.55) 的 $S_{IR,\mathcal{U}_{\Lambda^{2}}}$。
\item \textbf{数学工具：} 把扭曲后的 (5.55) 放到一般定向紧黎曼四维流形上，其 $Q$-不动点方程恰是 SW 方程 (5.57)（规范场是库仑分支上所用规范场的电磁对偶）。
\item \textbf{核心结论：} 一般四维流形上出现新的“量子希格斯分支”（即 SW 方程解的模空间），故 $Z^{\xi}_{IR}(p,s)=Z_{\text{Coulomb}}+Z_{\text{Higgs}}$ (5.58)；而 $Z_{\text{Higgs}}$ 本身是 $u=\Lambda^{2}$ 与 $u=-\Lambda^{2}$ 两个路径积分之和。
\end{keytakeaways}

既然已经回顾了 $\RR^{4}$ 上理论的 LEET，我们接下来考虑：把拓扑扭曲的理论表述在一般的连通、紧、定向黎曼四维流形 $X$ 上时会发生什么。

如上面提到的，由于 $X$ 体积有限，我们必须对所有真空积分。为说明这一点，应当考虑 $Y\times\RR$ 或 $Y\times S^{1}$ 上的路径积分。关键问题是不同真空之间是否存在隧穿。当 $Y$ 的体积无穷时隧穿振幅为零，我们在定义路径积分时取定无穷远处的某个真空。\footnote{〔原注〕这里重要的是我们在 $d>3$ 维时空的场论中工作；在更低维数，红外发散的问题要求对这一陈述作修改。}另一方面，当 $Y$ 的体积有限时，非零的隧穿振幅会把所有真空混合起来，因此必须在 $Y\times S^{1}$ 上的路径积分中对它们全部求和。对紧流形 $X$ 唯一自然的推广，就是继续对所有真空位形求和。对真空的求和有一部分是通过计算描述 $u$ 平面上涨落的 LEET 的配分函数来实现的。

然而在 $u=\pm\Lambda^{2}$ 的小邻域内，我们必须改用 (5.55) 给出的 LEET $S_{IR,\mathcal{U}_{\Lambda^{2}}}$ 或其在 $\mathcal{U}_{-\Lambda^{2}}$ 中的对应物。

作拓扑扭曲后，这又取 CohTFT 的标准形式。\footnote{〔原注〕前提是使用超多重态恰当的离壳表述；细节见 \cite{ref24}。}因此我们应研究这一组合理论的 $Q$-不动点方程。这些方程不是别的，正是我们上面讨论过的著名塞伯格--威滕方程 (4.74)：
\begin{equation}
\begin{aligned}
F(A_{d})^{+}&=\mu(M)\\
\not{D}M&=0
\end{aligned}
\tag{5.57}
\end{equation}

我们必须强调，(5.57) 中所用的规范场与库仑分支上其他地方所用的规范场是\emph{电磁对偶}的。

于是，在一般四维流形上\emph{存在一个新的量子“真空”分支}，我们称之为量子希格斯分支：它就是塞伯格--威滕方程解的模空间。\footnote{〔原注〕这些并不是通常意义上的真空：它们不是 $\RR^{4}$ 上理论的量子真空。但路径积分显然必须局域化到这一轨迹上。}因此，$Z^{\xi}_{DW}(p,s)$ 的红外计算涉及两项之和：
\begin{equation}
\boxed{Z^{\xi}_{IR}(p,s)=Z_{\text{Coulomb}}+Z_{\text{Higgs}}}
\tag{5.58}
\end{equation}
当然，$Z_{\text{Higgs}}$ 这一项本身又是两个路径积分之和：一个给出 $u=\Lambda^{2}$ 处的贡献，另一个给出 $u=-\Lambda^{2}$ 处的贡献。我们常把 $Z_{\text{Coulomb}}$ 简记为 $Z_{u}$，因为正如我们将看到的，它化为 $u$ 平面上的一个积分。

现在，为了显式算出 $Z_{\text{Coulomb}}$ 与 $Z_{\text{Higgs}}$，我们需要补全一般黎曼流形 $X$ 上理论作用量的描述。塞伯格与威滕确定的是他们在 $\RR^{4}$ 上的有效作用量；当耦合到引力场时，会出现必须计入的新项。

\subsubsection{库仑分支上的引力耦合}

先讨论库仑分支上有效作用量中的耦合。由于拓扑不变性，与度规的耦合只能通过
\begin{equation}
\Delta_{\text{grv}}S=\int_{X}e(u)\Tr R\wedge R^{*}+p(u)\Tr R\wedge R
+\frac{\ii}{4}\int_{X}F\wedge w_{2}(X)
\tag{5.59}
\end{equation}
出现，其中 $F$ 是低能 $U(1)$ 阿贝尔规范理论的场强。

前两项在 $u$ 平面上的积分中指数化为函数
\begin{equation}
E(u)^{\chi}P(u)^{\sigma}
\tag{5.60}
\end{equation}
现在，由物理论证可以导出：
\begin{equation}
E(u)=\alpha\left(\frac{\dd u}{\dd a}\right)^{1/2}
\tag{5.61}
\end{equation}
\begin{equation}
P(u)=\beta\Delta^{1/8}
\tag{5.62}
\end{equation}
其中 $\alpha,\beta$ 为常数（可以依赖于 $\Lambda$，但不依赖于 $u$）。

% [Tutor Note: (5.61)(5.62) 的“物理论证”作者一笔带过，只说“one can derive from physical arguments”。其实质是三条约束的交汇：(i) 拓扑不变性只允许 (5.59) 那种与度规的耦合，而 $R\wedge R^{*}$、$R\wedge R$ 的系数在四维流形上分别按 $\chi$ 与 $\sigma$ 指数化，故积分中只能出现 $E^{\chi}P^{\sigma}$ 的组合；(ii) $E(u)$、$P(u)$ 必须对 $u$ 全纯——理由是它们来自只含整数量子数的费米子测度的单值化：把理论放在 $S^{1}\times M_{3}$ 上绕一圈，$R$-对称性反常给出测度在 $u$ 平面上的单值性，而 $\tau$ 的单值变换同时作用在 $E$、$P$ 上，唯一能与 $\mathrm{Im}\,\tau\to\infty$ 处的弱耦合结果相容的全纯解就是 (5.61)(5.62) 的形式（$\Delta^{1/8}$ 中的 $1/8$ 正是由 $\Delta$ 在奇异点附近按 $\delta u^{2}$ 消失、而 $P$ 需按 $\delta u^{1/4}$ 消失定出的）；(iii) $\alpha,\beta$ 是纯数，不含 $u$，因此只能由“外部数据”标定：要么用两类特殊流形上唐纳森多项式的显式计算，要么用穿墙公式，要么用爆破公式——三者给出同一答案 (5.65)。作者在 §5.7.1 末尾正是选择了穿墙公式这条路线。]

\medskip
\noindent\textbf{评注：}
\begin{enumerate}[label=\arabic*., leftmargin=2.1em]
\item 这里 $\Delta$ 是曲线的判别式：
\begin{equation}
\Delta\sim\prod_{i}(u-u_{i})
\tag{5.63}
\end{equation}
其中 $u_{i}$ 是 BPS 态变为无质量的位置。\footnote{〔原注〕正如 \cite{ref49} 所强调的，一般而言这\emph{不是}描述秩一理论的椭圆曲线的判别式。〔原稿边注〕需要澄清它是否是谱曲线形式的分支轨迹——我认为应当是。}对纯 SYM，它被归一化得使 $u\to\infty$ 时\footnote{〔原稿边注〕幂次 $\kappa$ 由量纲分析确定。}
\begin{equation}
\Delta=\Lambda^{\kappa}\prod_{\alpha}Z_{\alpha}^{2}
\tag{5.64}
\end{equation}
\item \cite{ref38,ref56} 中给出的推导依赖于 $u$ 平面上被积函数的一致性、$u\to\infty$ 时的已知行为，特别是它作为 $R$-对称性变换下的函数的行为。
\item 我们把 $E$ 与 $P$ 写成了一种在把 $SU(2)$ vm 耦合到超多重态时仍可推广的形式（这在后文讨论 geography 问题时会很重要）。在纯 $SU(2)$ 理论中 $\Delta=u^{2}-\Lambda^{4}$。
\item 这里 $\alpha,\beta$ 是数值常数。最终它们将是整个计算中\emph{唯一}的未知量，并将被“拟合到实验数据”。它们可由以下方式提供：
\begin{enumerate}[label=\alph*.), leftmargin=1.8em]
\item 两类特殊情形下唐纳森多项式的显式计算；
\item 穿墙公式；
\item 爆破公式。
\end{enumerate}
当然，三种方法给出同一答案。在本书中我们将使用方法 (b)——与数学上导出的穿墙公式相匹配。结果为\footnote{〔原稿边注〕与拉巴斯蒂达--马里诺式 8.17 不同，为什么？}
\begin{equation}
\alpha=\left(\frac{2}{\pi^{2}}\right)^{1/8}\qquad
\beta=\left(\frac{27}{\pi^{2}}\right)^{1/8}
\tag{5.65}
\end{equation}
\end{enumerate}

% [Tutor Note: (5.63)(5.64) 的 $\Delta$ 与 §5.2 的 $\Delta$ 是两回事：前者是“BPS 态变轻的位置”的判别式（重根轨迹），后者是库仑分支上“半单但非正则”的超平面并集。作者在此用同一字母，且脚注 28 又强调它一般不是椭圆曲线的判别式——这一点在把结论推广到任意库仑分支时至关重要：对秩一理论二者偶然吻合（纯 SYM 时 $\Delta=u^{2}-\Lambda^{4}$），但一般情形下应理解为谱曲线的分支轨迹。另注意 (5.64) 中 $\prod_\alpha Z_\alpha^{2}$ 的 $Z_\alpha$ 是根 $\alpha$ 对应的“质量”因子，其平方给出 $\Delta$ 在弱耦合端的标度行为；幂次 $\kappa$ 由量纲分析定出（纯 SU(2) 时 $\Lambda^4$ 与 $u^2$ 同量纲，故 $\kappa=4$）。]

(5.59) 中的第三项是 LEET $U(1)$ VM 中阿贝尔规范场 $A$ 在通量扇区求和中出现的一个重要相位：
\begin{equation}
Z_{\text{Coulomb}}=\sum_{\lambda}\langle\cdots\rangle_{[F]=4\pi\lambda}
\tag{5.66}
\end{equation}
为准确陈述这一相位，回忆在 UV 中我们有一个带 't Hooft 磁通 $\xi\in H^{2}(X,\ZZ/2\ZZ)$ 的 $SO(3)$ 丛 $P$。我们取 $SO(3)$ 丛 $P$ 的特征类 $\xi$ 的一个整提升 $\bar\xi$（假定这样的提升存在），把它嵌入 $H^{2}(X,\RR)$，并将其在 $H^{2}(X;\RR)$ 中的像除以二，以定义\footnote{〔原稿边注〕我们应当把记号从 $w_{2}(E)$ 改成 $w_{2}(P)$，其中 $P$ 是 $SO(3)$ 丛。这样更有意义，何况 $E$ 本来就用得太滥了。}
\begin{equation}
2\lambda_{0}=\bar\xi
\tag{5.67}
\end{equation}
于是低能阿贝尔规范理论中出现的线丛，其“第一陈类”落在挠子
\begin{equation}
\lambda\in\Gamma_{\xi}:=\lambda_{0}+\bar H^{2}(X)
\tag{5.68}
\end{equation}
中。\footnote{〔原注〕这里我们又一次对挠不严谨。}

理解 $\lambda_{0}$ 的一种方式如下：当我们有一个 $SU(2)$ 规范丛 $P$ 时，利用旋量表示有与之相伴的秩二复向量丛 $Q$。自发对称破缺把规范对称性破缺到 $U(1)$ 后，$Q\cong L\oplus L^{-1}$，其中 $L$ 是线丛。对瞬子扇区的求和就变成对线丛 $L$ 的拓扑类求和。然而，当 $P$ 是具有非平凡 $w_{2}(P)$（即非平凡 't Hooft 磁通）因而不能提升为 $SU(2)$ 丛的 $SO(3)$ 丛时，$Q$ 与 $L$ 并不存在。它们存在性的障碍可以从那些破坏三重坐标重叠上链条件余循环的 $\pm1$ 中看出。这些障碍在考虑对称平方
\begin{equation}
\Sym^{2}(L\oplus L^{-1})=L^{2}\oplus\OO\oplus L^{-2}
\tag{5.69}
\end{equation}
时消失；这里我们记 $2\lambda=c_{1}(L^{2})$，于是
\begin{equation}
c_{1}(L^{2})=2\lambda=\bar\xi \mod 2\bar H^{2}(X)
\tag{5.70}
\end{equation}
计算 $Z_{\text{Coulomb}}$ 时需要对这样的线丛求和，因此路径积分是对 $\bar H^{2}(X)$ 的一个挠子的求和。

在 \cite{ref56} 中，威滕论证了：对为得到 LEET 而被积分掉的\emph{有质量}模式，路径积分测度存在一个相对相位。唯一的问题在于带电费米子测度的定向。回忆 $\psi,\bar\psi$（局部地）取值于 $S^{\pm}\otimes S_{R}\otimes\mathfrak{g}$，而其中取值于 $\mathfrak{g}-\mathfrak{t}$ 的部分在未破缺的、李代数为 $\mathfrak{t}$ 的规范对称性下带电。

% [Tutor Note: 通量扇区求和中的相对相位来自三处，容易混淆。(i) (5.66) 的求和变量是 $\lambda=\tfrac12 c_{1}(L^{2})$，即“第一陈类的一半”，落在挠子 $\Gamma_{\xi}$ 上，而非普通上同调类。(ii) (5.59) 的第三项 $\tfrac{i}{4}\int F\wedge w_{2}(X)$ 在扇区 $[F]=4\pi\lambda$ 上取值 $\ee^{\,\ii\pi\lambda\cdot w_{2}(X)}$，这是“电”相位，随 $\lambda$ 线性；它正是 (5.59) 第三项存在的理由。(iii) 积分掉有质量带电费米子时，测度定向给出相对符号 (5.71)，它只依赖差 $\lambda_{1}-\lambda_{2}$，即 $(-1)^{w_{2}(X)\cdot(\lambda_{1}-\lambda_{2})}$。把 (ii) 与 (iii) 合起来、再补上 $\lambda_{0}$ 处的二次项 $\ee^{\pm2\pi i\lambda_{0}^{2}}$，就得到总相位 (5.72)，其中 $\lambda_{0}^{2}$ 为自交数；取 $+$ 号等价于选取模空间 $\mathcal{M}(P,\mathfrak{g})$ 的定向 (5.73)。(5.74) 的平移只把总相位改变 $(-1)^{\beta\cdot w_{2}(X)}=(-1)^{\beta^{2}}$，因为 $w_{2}(X)$ 是特征向量——这正是“改变整提升不改变物理”的具体体现，也是后文处理 $u$ 平面积分单值性的基础。]

人们可以试着用（比如说）一个殆复结构来为 UV 理论中的费米子测度选定标准定向。\footnote{〔原注〕殆复结构存在当且仅当存在 $c\in\bar H^{2}(X)$ 使 $c^{2}=2\chi+3\sigma$。我们将看到，若此条件不满足，则 DW 与 SW 理论都是空的。}在结构群约化之后，对每个 $\lambda$ 的选取，同一组费米子的路径积分测度都有一个不同的标准定向。威滕论证后者是正的，并用一个 Dirac 算子的指标计算了相对定向。最终结果是：积分掉有质量费米子模式对 $\lambda$ 的依赖，在扇区 $\lambda_{1}$ 与 $\lambda_{2}$ 之间给出相对符号
\begin{equation}
(-1)^{w_{2}(X)\cdot(\lambda_{1}-\lambda_{2})}
\tag{5.71}
\end{equation}
这仍然留下一个整体相位需要拟合。若取
\begin{equation}
\ee^{\pm2\pi\ii\lambda_{0}^{2}}(-1)^{(\lambda-\lambda_{0})\cdot w_{2}(X)}
\tag{5.72}
\end{equation}
则下文 $Z_{DW}$ 的计算（当 $\Lambda$ 为实时为实）。第一个指数中符号的选取对应于模空间 $\mathcal{M}(P,\mathfrak{g})$ 定向的选取。我们将取 $+$ 号：
\begin{equation}
\ee^{2\pi\ii\lambda_{0}^{2}}(-1)^{(\lambda-\lambda_{0})\cdot w_{2}(X)}
\tag{5.73}
\end{equation}
最后，为后文引用，注意若平移原点
\begin{equation}
\lambda_{0}\to\lambda_{0}+\beta
\tag{5.74}
\end{equation}
即把提升平移为 $\bar\xi\to\bar\xi+2\beta$，则总相位改变
\begin{equation}
(-1)^{\beta\cdot w_{2}(X)}=(-1)^{\beta^{2}}
\tag{5.75}
\end{equation}
因为 $w_{2}(X)$ 是特征向量。

\medskip
\noindent\textbf{评注：}
\begin{enumerate}[label=\arabic*., leftmargin=2.1em]
\item 上述讨论与唐纳森--克龙海默 \cite{ref8} 第 7.1.6 节的讨论密切相关，尤其见第 283 页最后一段。
\end{enumerate}

\subsubsection{希格斯分支上的引力耦合}

希格斯分支上的引力耦合有类似的故事。(5.55) 的作用量在拓扑扭曲后取如下形式
\begin{equation}
Q(\Psi)+\int_{X}c(u)F_{d}\wedge F_{d}
\tag{5.76}
\end{equation}
其中 $F_{d}=F(A_{d})$ 是\emph{对偶光子}的曲率。引力耦合将具有如下形式：
\begin{equation}
\Delta_{\text{grv}}S_{Higgs}=\int_{X}e_{h}(u)\Tr R\wedge R^{*}
+p_{h}(u)\Tr R\wedge R+\frac{\ii}{4}F_{d}\bar\xi
\tag{5.77}
\end{equation}
这四项中，最后一项可以由 (5.59) 中的相应项\emph{导出}。这来自从适用于远离 $u=\pm\Lambda^{2}$ 的 $u$ 平面的对偶标架，变到适用于磁单极点与双荷子点的对偶标架所必需的电磁对偶变换。细节最好留到我们写出 $u$ 平面被积函数之后再讲，但结论将是：光子量子化与耦合由下式给出。

\noindent\underline{远离 $u=\pm\Lambda^{2}$：}光子的量子化由
\begin{equation}
  [F]=4\pi\lambda ,
  \tag{5.78}
\end{equation}
给出，其中 $\lambda\in\tfrac{1}{2}\bar\xi+\bar H^{2}(X)$。相互作用为
\begin{equation}
\ee^{\int_{X}\frac{\ii}{4}F\wedge w_{2}(X)}=\ee^{\ii\pi\lambda\cdot w_{2}(X)}.
\tag{5.79}
\end{equation}

\noindent\underline{在 $u=\pm\Lambda^{2}$ 附近：}光子的量子化由
\begin{equation}
  [F_{d}]=4\pi\lambda ,
  \tag{5.80}
\end{equation}
给出，其中 $\lambda\in\tfrac{1}{2}w_{2}(X)+\bar H^{2}(X)$。相互作用为
\begin{equation}
\ee^{\int_{X}\frac{\ii}{4}F_{d}\wedge\bar\xi}=\ee^{\ii\pi\lambda\cdot\bar\xi}
\tag{5.81}
\end{equation}
特别要注意，在式 (5.80) 中类 $2\lambda$ 是某个 Spin$^{c}$ 结构的特征类。

其余各项的性质则很不一样。它们指数化为
\begin{equation}
\ee^{2\pi\ii\left(\lambda_{0}^{2}+\lambda\cdot\lambda_{0}\right)}
C(u)^{\lambda^{2}}P_{h}(u)^{\sigma}E_{h}(u)^{\chi}
\tag{5.82}
\end{equation}
其中现在 $\lambda$ 是一个 spin$^{c}$ 结构，而 $2\lambda_{0}$ 是我们先前用过的 $w_{2}(P)$ 的整提升。关键之点在于：这三个函数 $C,P,E$ 至今尚未确定，且不易从第一性原理导出。

然而，函数 $C,P,E$ 是普适的——意思是它们与四维流形 $X$ 无关——这一点连同 $u$ 平面积分的穿墙现象，将使我们能够显式地确定它们。\footnote{〔原稿边注〕其实 $C(u)$ 应当很容易由积分掉单极子场并与库仑分支作用量相匹配而导出！}\footnote{〔原稿边注〕距 §5.7.2 开始已过 6 分钟。}

\subsection{把算子从 UV 映射到 IR}

\begin{keytakeaways}
\item \textbf{物理图像：} 按定义 $\mathcal{O}$ 就是 $2u$，这在 UV 与 IR 都成立；但用场来表达 $u$ 的式子在两个理论中差别很大，因此必须把 UV 算子译成 IR 的场。
\item \textbf{数学工具：} 诀窍在于 (4.36) 的 $K$ 算符是 RG 不变的（超对称算符不随标度演化）：只需用 IR 的场与超对称变换律把 $K$ 作用两次于 $2u$，即得 (5.84)（库仑分支）与 (5.85)（希格斯分支，此时 $u$ 是 $a_{d}$ 的已知模函数）。
\item \textbf{核心结论：} RG 映射不是环同态：面算子相乘时会多出正比于相交数 $\Sigma_{1}\cdot\Sigma_{2}$ 的接触项 $\Sigma^{2}T(u)$ (5.87)，它使关联函数指数化为 (5.88)(5.89)，并给出 $T_{c}(u)$ 的模形式表达式 (5.91) 及 $\hat E_{2}$ 的修正 (5.92)。
\end{keytakeaways}

现在必须理解如何把算子 $\mathcal{O}$ 与 $\mathcal{O}(\Sigma)$ 表达为低能有效理论中的算子。

诀窍在于理解 0-算符在 RG 下如何映射，然后认识到 (4.36) 的 $K$ 算符是 RG 不变的：超对称算符不随标度演化。

\subsubsection{库仑分支上的算子映射}

按定义，算子 $\mathcal{O}$ 就是 $2u$。这在低能与高能都成立。用场来表达 $u$ 的式子在 UV 与 IR 理论中会很不一样。

无论如何，在 IR 理论中，我们用在低能有效阿贝尔理论中的场与超对称变换律，把 $K$ 作用于 $2u$ 上而得到 $\mathcal{O}^{(1)}$。于是利用标准的超对称变换可得
\begin{equation}
\begin{aligned}
Ka&\sim\psi\\
K\psi&\sim(F^{-}+D)
\end{aligned}
\tag{5.83}
\end{equation}
等等。（回忆 $D$ 是自对偶 2-形式。）

于是在低能理论中 $\mathcal{O}^{(1)}=Ku\sim\frac{\partial u}{\partial a}\psi$；再作用一次 $K$，得到
\begin{equation}
\begin{aligned}
\mathcal{O}_{IR,c}&=2u\\
\mathcal{O}_{IR,c}(\Sigma)&\sim\int_{\Sigma}\frac{\partial u}{\partial a}\left(F^{-}+D\right)
+\frac{\partial^{2}u}{\partial a^{2}}\psi^{2}
\end{aligned}
\tag{5.84}
\end{equation}

\subsubsection{希格斯分支上的算子映射}

类似地，在希格斯分支上只有一个具有正确鬼数的 0-算符，它就是 $a_{d}$。算子 $\mathcal{O}=u$ 是 $a_{d}$ 的已知函数，用模函数表出。因此，完全按照我们在库仑分支上所用的策略，得到低能算子
\begin{equation}
\begin{aligned}
\mathcal{O}_{IR,h}&=2u\\
\mathcal{O}_{IR,h}(\Sigma)&\sim\int_{\Sigma}\frac{\dd u}{\dd a_{d}}F(A_{d})
+\frac{\dd^{2}u}{\dd a_{d}^{2}}\psi^{2}
\end{aligned}
\tag{5.85}
\end{equation}
其中 $A_{d}$ 是在单极子纯粹带电荷的对偶标架中的 $U(1)$ 规范场。

\subsubsection{接触项}

在低能有效理论中计算算子 $\mathcal{O}(\Sigma)$ 的关联函数时，有一个重要的微妙之处。当 $\Sigma$ 有自交点时，即使在拓扑场论中也会出现必须计入的奇异性。人们考虑的是路径积分中算子的缩并：
\begin{equation}
\left\langle\cdots\int_{\Sigma_{1}}\Tr\left(\phi(x_{1})F_{\mu\nu}+\cdots\right)\dd x_{1}^{\mu}\dd x_{1}^{\nu}
\int_{\Sigma_{2}}\Tr\left(\phi(x_{2})F_{\lambda\rho}+\cdots\right)\dd x_{2}^{\mu}\dd x_{2}^{\nu}\cdots\right\rangle
\tag{5.86}
\end{equation}

由于拓扑不变性，来自 $x_{1}\neq x_{2}$ 的场缩并的贡献将合并为全导数，结果只依赖于 $\Sigma_{1}$ 与 $\Sigma_{2}$ 的拓扑。然而，在 $\Sigma_{1}$ 与 $\Sigma_{2}$ 的交点 $x_{1}=x_{2}$ 处，这些全导数会给出贡献。再次援引拓扑不变性（这里使用离壳超对称），它必须是某个与点相联系的 $Q$-闭算子，因而只是 $u$ 的全纯函数；结果必定可以用 $\wp=x_{1}=x_{2}$ 处某个局域 $Q$-不变算子来表达。

\begin{figure}[H]
  \centering
  \figplaceholder{〔原文此处于第 57 页留有图位，PDF 中未插入实际图形〕}
  \caption{两个相交的曲面。}
  \label{fig:intersectingsurfaces}
\end{figure}
% [原 PDF 第 57 页]

换个说法：RG 映射 $\mathcal{O}_{UV}\to\mathcal{O}_{IR}$ 在面算子上\emph{不是}环同态。确切地说，
\begin{equation}
\mathcal{O}_{UV}(\Sigma_{1})\mathcal{O}_{UV}(\Sigma_{2})
\to\mathcal{O}_{UV}(\Sigma_{1})\mathcal{O}_{UV}(\Sigma_{2})
+\Sigma_{1}\cdot\Sigma_{2}T(u)+Q(\ast)
\tag{5.87}
\end{equation}
其中 $T(u)$ 是某个 0-可观测（量）。注意我们这里用到了 $\mathcal{O}(\wp_{1})-\mathcal{O}(\wp_{2})=Q(\ast)$。由于 0-可观测之间没有这类奇异性，我们有如下的 Wick 定理：
\begin{equation}
\left\langle\ee^{\mathcal{O}_{UV}(\Sigma)}\right\rangle_{UV}
=\left\langle\ee^{\mathcal{O}_{IR}(\Sigma)+\Sigma^{2}T(u)}\right\rangle_{IR}
\tag{5.88}
\end{equation}
原则上，算子 $T(u)$ 可以依赖于真空分支，因此现在我们有
\begin{equation}
\begin{aligned}
Z^{\xi}_{DW}(p,\Sigma)&=\left\langle\ee^{p\mathcal{O}+\mathcal{O}(\Sigma)}\right\rangle_{UV}\\
&=\left\langle\ee^{p\mathcal{O}_{IR,c}+\mathcal{O}_{IR,c}(\Sigma)+\Sigma^{2}T_{c}(u)}\right\rangle_{\text{Coulomb}}
+\left\langle\ee^{p\mathcal{O}_{IR,h}+\mathcal{O}_{IR,h}(\Sigma)+\Sigma^{2}T_{h}(u)}\right\rangle_{\text{Higgs}}
\end{aligned}
\tag{5.89}
\end{equation}
函数 $T_{c}(u)$ 与 $T_{h}(u)$ 可由自洽性论证确定，\cite{ref38} 中就是这样做的。\footnote{〔原稿边注〕也许干脆把 $T(u)$ 留作另一个未知函数，待我们讨论如何使 $u$ 平面积分良定义时再定。}这些接触项的系统理论由洛谢夫--涅克拉索夫--沙塔什维利 \cite{ref26} 发展起来；结果的简洁推导另见 \cite{ref29}：
\begin{equation}
T_{c}(u)\sim\frac{\partial^{2}\mathcal{F}}{\partial\tau_{0}^{2}}
\qquad \tau_{0}\sim\log\Lambda.
\tag{5.90}
\end{equation}
对纯 $SU(2)$ 理论，这是对偶群 $\Gamma^{0}(4)$ 下的某个权零几乎模形式：
\begin{equation}
T_{c}(u)=-\frac{1}{24}\left(\left(\frac{\dd u}{\dd a}\right)^{2}E_{2}-8u\right)
\tag{5.91}
\end{equation}
我们还用 $\hat T_{c}$ 表示把 $E_{2}$ 换成标准的模（但非全纯）对象
\begin{equation}
\hat E_{2}(\tau)=E_{2}(\tau)-\frac{3}{\pi y}
\tag{5.92}
\end{equation}
所得的量；这里以及下文，我们记 $\tau=x+\ii y$ 为 $\tau$ 的实部与虚部。\footnote{〔原稿边注〕距 §5.8 开始已过 10 分钟。}

% [Tutor Note: (5.90) 只是示意性的：$T_{c}(u)\sim\partial^{2}\mathcal{F}/\partial\tau_{0}^{2}$，配合 $\tau_{0}\sim\log\Lambda$ 表明接触项量度的是对 UV 耦合的依赖，量纲上每次 $\partial_{\tau_{0}}$ 带出 $\Lambda^{2}$ 因子，故 $\Sigma^{2}T_{c}$ 与 $\mathcal{O}(\Sigma)$ 的量纲相容。真正可操作的是 (5.91)：它是 $\Gamma^{0}(4)$ 下权零的“几乎模形式”。$E_{2}$ 本身只按权 $2$ 的模变换“差一个非全纯项”变换（见 (5.92)），故单独 $(du/da)^{2}E_{2}$ 不变；减去 $-8u$ 后整体成为权零不变对象，其中 $u$ 是权零模函数。把 $E_{2}$ 换成 $\hat E_{2}$ (5.92) 后，$\hat E_{2}$ 严格按权 $2$ 变换，从而 $u$ 平面积分的被积函数在 $\Gamma^{0}(4)$ 下良定义；这也解释了为何后文只需处理 $\hat T_{c}$。注意 (5.90) 中 $\tau_{0}$ 是裸耦合 (5.17)，不是 LEET 的 $\tau(a)$，二者不可混淆。]

% [Tutor Note: 本文件到此为止，正文原书接下来是 §6“General Form Of The Higgs Branch Contribution”。原文 §5 的图位共 5 处，均已用 figure 环境复现。]

